\originalchapter{作业五：正确的玻尔兹曼计数}
\psource{8.044，2013 春，PS5，3 月 11 日截止。原卷只有一道题，含 (a)–(c)；据官方答案核对的 AI 独立解答。}
\originalsection{1．广延性与 Gibbs 悖论}
\begin{exercise}
未校正的单原子理想气体计数为 $\Phi=V^N(4\pi e m E/(3N))^{3N/2}$。
(a) 用 $S=\kb\ln\Phi$、$E=3N\kb T/2$ 表示 $S(N,V,T)$；把所有广延量乘 $\lambda$、强度量不变，检验是否成为 $\lambda S$。
(b) 可动隔板分出 $V_1=\alpha V,V_2=(1-\alpha)V$，两种气体初始同温同压，证明 $N_1=\alpha N,N_2=(1-\alpha)N$；撤板混合求各气体及总熵增。若两边是同种气体，说明 $\Delta E=\Delta W=\Delta Q=0$，可逆撤插隔板为何要求 $\Delta S=0$，与原计数比较。
(c) 同种粒子的置换不产生新微态，将 $\Phi,\Omega,D$ 除以 $N!$，直接证明解决 (a)(b) 两问题；可用大 $N$ 的 Stirling 近似。
\end{exercise}
\begin{solution}
(a) 加入固定态单元 $h^{3N}$ 使对数无量纲，$S_{\rm old}=N\kb\ln[V(2\pi em\kb T/h^2)^{3/2}]$，缩放后多出 $\lambda N\kb\ln\lambda$，非广延。
(b) 同温同压给 $N_i/V_i=P/(\kb T)=N/V$。不同种气体各粒子数不变，$\Delta S_1=\alpha N\kb\ln(1/\alpha)$，$\Delta S_2=(1-\alpha)N\kb\ln[1/(1-\alpha)]$，总和非负。理想气体总能仅依赖 $N,T$，同种撤板不改变二者，且无边界移动，故能、功、热皆零。两侧本来具有相同宏观状态，隔板可可逆恢复；若计数仍给上述正混合熵即为 Gibbs 悖论。
(c) $S=S_{\rm old}-\kb\ln N!\simeq N\kb[\ln(V/(N\lambda_T^3))+5/2]$，$\lambda_T=h/\sqrt{2\pi m\kb T}$，缩放 $V,N$ 后括号不变。不同种分别除 $N_1!,N_2!$，混合熵不变；同种则撤板后以一个 $N!$ 计数，初态同密度的两块熵相加恰等于末态熵。有限尺寸隔板位置约束可有次广延涨落修正，不改变热力学极限结论。原题允许略去 $O(N)$ 项，但保留 Stirling 的 $-N$ 更便于得到标准熵常数。
\end{solution}
